\section{复位状态、ROM 映射与首条指令}

\topic{复位是契约，不是初始化完成}
上电复位后，处理器处于实地址模式。最关键的可观察状态是指令入口：
\code{CS} 的可见值为 \code{0xF000}，隐藏基址为
\code{0xFFFF0000}，\code{IP/EIP} 指向 \code{0xFFF0}，因此第一条
指令的物理地址是：
\[
  0xFFFF0000 + 0xFFF0 = 0xFFFFFFF0.
\]
隐藏基址是理解现代 x86 复位地址的关键。此时不能简单套用普通实模式
\(segment \times 16 + offset\) 去解释第一次取指。

\begin{pdfdiagram}
  \node[os hardware] (reset) {RESET};
  \node[os state,below left=of reset] (visible) {CS 可见值\\\code{0xF000}};
  \node[os block,below=of reset] (hidden) {CS 隐藏基址\\\code{0xFFFF0000}};
  \node[os block,below right=of reset] (offset) {EIP\\\code{0x0000FFF0}};
  \node[os hardware,below=of hidden] (physical) {首条物理地址\\\code{0xFFFFFFF0}};
  \node[os state,below=of physical] (reload) {far transfer 重载 CS\\恢复
  selector \(\times16\) 规则};
  \draw[os arrow] (reset) -- (visible);
  \draw[os arrow] (reset) -- (hidden);
  \draw[os arrow] (reset) -- (offset);
  \draw[os arrow] (hidden) -- node[above left]{相加} (physical);
  \draw[os arrow] (offset) -- (physical);
  \draw[os arrow] (physical) -- (reload);
\end{pdfdiagram}
\diagramnote{依据 Intel SDM Vol. 3A 12.1.4 重绘。手册给出复位时 CS
可见值、隐藏基址和 EIP；项目再把结果映射到 128 KiB ROM 的文件偏移。}

\topic{128 KiB ROM 的顶端映射}
项目 ROM 覆盖 \code{0xFFFE0000} 到 \code{0xFFFFFFFF}。文件偏移与物理
地址之间满足：
\[
  physical = 0xFFFE0000 + file\_offset.
\]
所以复位向量的文件偏移是 \code{0x1FFF0}，入口
\code{0xFFFFF000} 的文件偏移是 \code{0x1F000}。

\begin{center}
\begin{tikzpicture}[os diagram]
  \node[os hardware,minimum width=42mm] (rom) {ROM 起点\\0xFFFE0000};
  \node[os state,above=12mm of rom,minimum width=42mm] (entry)
    {16 位入口\\0xFFFFF000};
  \node[os block,above=12mm of entry,minimum width=42mm] (reset)
    {复位向量\\0xFFFFFFF0};
  \node[os hardware,above=8mm of reset,minimum width=42mm] (end)
    {ROM 末端\\0xFFFFFFFF};
  \draw[os arrow] (rom) -- node[right]{128 KiB 地址空间} (entry);
  \draw[os arrow] (entry) -- (reset);
  \draw[os arrow] (reset) -- (end);
\end{tikzpicture}
\end{center}

\topic{16 位近跳转保留复位隐藏状态}
复位向量放置操作码 \code{E9} 和有符号 16 位位移。近跳转只修改 IP，不
重载 CS，因而保留隐藏基址 \code{0xFFFF0000}。目标 IP 是
\code{0xF000}，物理目标正好是 \code{0xFFFFF000}。位移为：
\[
  0xF000 - (0xFFF0 + 3) = -4083 = 0xF00D.
\]
因此复位向量前三个字节是 \code{E9 0D F0}。

\begin{failurebox}
如果使用会重载 CS 的远跳转，就必须重新证明目标段如何映射到 ROM；如果链接
脚本把入口放错一个字节，CPU 也不会给出友好错误，只会执行填充值或触发异常。
\end{failurebox}

\topic{复位向量固定在末端的历史连续性}
8086 只有 20 位物理地址总线，复位后从物理地址 \code{0xFFFF0} 开始执行，
这个位置距 1 MiB 地址空间末端只有 16 字节。把入口放在顶端，使主板能够把
固件映射在高地址，同时把大部分低地址留给 RAM。后续处理器扩展了地址宽度，
但软件生态已经依赖这一入口语义。现代 x86 在复位时借助 CS 的可见选择子和
隐藏基址组合，把第一条取指定位到 \code{0xFFFFFFF0}，继续保持“距 4 GiB
末端 16 字节”的结构。

兼容性使新处理器可以运行既有固件，也使现代启动过程必须从一个极受约束的
历史状态开始。AMD64 没有把复位入口直接改成长模式，因为这会同时破坏固件
生态、主板设计和诊断工具。长模式因此是软件逐步建立的目标状态，而不是上电
默认状态。

\topic{近跳转字节的逐项推导}
复位位置只有 16 字节可供使用，真正的初始化代码必须放在 ROM 的其他位置。
\code{JMP rel16} 的机器编码由一个字节操作码和一个小端序的 16 位有符号
位移组成。处理器先把 IP 推进到下一条指令，再把位移加到 IP：
\[
  \mathrm{IP}_{next}=0xFFF0+3=0xFFF3,
  \qquad
  \mathrm{IP}_{target}=0xFFF3+\operatorname{signext}(0xF00D)=0xF000.
\]

\begin{osgridlongtable}{L{0.16\textwidth}L{0.22\textwidth}L{0.48\textwidth}}
\textbf{字节位置} & \textbf{数值} & \textbf{处理器解释}\\ \hline
\endfirsthead
\textbf{字节位置} & \textbf{数值} & \textbf{处理器解释}\\ \hline
\endhead
\code{0xFFFFFFF0} & \code{E9} & 选择近相对跳转指令\\ \hline
\code{0xFFFFFFF1} & \code{0D} & 位移低字节\\ \hline
\code{0xFFFFFFF2} & \code{F0} & 位移高字节；组合为 \code{0xF00D}\\ \hline
\code{0xFFFFFFF3} & \code{FF} & 不属于该指令；正常执行不会到达\\ \hline
\end{osgridlongtable}

近跳转不写 CS，因此隐藏基址保持 \code{0xFFFF0000}。新的线性地址为
\code{0xFFFF0000 + 0xF000 = 0xFFFFF000}，正好落在链接器安排的 ROM
入口。这个推导同时连接了机器码、小端序、符号扩展、指令长度、段缓存和链接
布局，任何一个量变化都必须重新证明。

\topic{填充值也是系统行为}
ROM 未占用区域填充为 \code{0xFF}，能够使镜像大小和烧录语义保持确定，但
\code{0xFF} 并不是安全的“空指令”。控制流意外进入填充区后，处理器可能将
连续字节解释为有效指令前缀与操作码，也可能触发异常。早期尚未建立异常处理，
最终表现通常只是复位、停机或无输出。因此链接断言和二进制字节检查必须在
运行前阻止布局错误。
